\chapter{Maxwell's Equations, Covariance, and Electromagnetic Waves}

\lectureoverview{Electric and magnetic fields act on charged matter through the Lorentz force, but matter must also create and rearrange the fields. We shall formulate that reciprocal half of electromagnetism, show that Maxwell's equations are Lorentz covariant, derive their wave solutions, and then restore charge and current in a form that makes local conservation manifest. Throughout the lecture we use units in which \(c=1\).}

\section{The Other Half of Electromagnetism}

Begin with the part of the subject already in hand. If a particle of charge \(q\) moves with velocity \(\mathbf v\) through electric and magnetic fields, the force is
\begin{equation}
\mathbf F=q\bigl(\mathbf E+\mathbf v\times\mathbf B\bigr).
\label{eq:lorentz-force-three}
\end{equation}
The factor \(q\) is indispensable. An uncharged particle does not feel the electromagnetic force. The electric field contributes even when the particle is instantaneously at rest, whereas the magnetic contribution depends on its motion.

Equation~\eqref{eq:lorentz-force-three} tells us what prescribed fields do to matter. It does not tell us what determines the fields. What produces an electric field around one object, a magnetic field around another, or a wave moving through empty space? The equations that answer those questions are Maxwell's equations.

\begin{figure}[H]
\centering
\includegraphics[width=0.62\textwidth]{lecture_58_figure_02.png}
\caption{The three-vector Lorentz force and the beginning of its tensor reformulation.}
\lectureframe{00:05:49}
\end{figure}

The relativistic form of the force law anticipates the notation we shall need:
\begin{equation}
f^\mu=qF^\mu{}_{\nu}u^\nu.
\label{eq:lorentz-force-four}
\end{equation}
Here \(u^\nu\) is the four-velocity, \(f^\mu\) is the four-force, and \(F^{\mu\nu}\) is an antisymmetric tensor assembled from \(\mathbf E\) and \(\mathbf B\). The board first displays only the schematic beginning of this expression; the contraction with \(u^\nu\) supplies the missing vector needed to leave one free force index.

There is a physical reason that matter must act back on the field. An electromagnetic wave can strike a charge, accelerate it, and transfer energy to it. If the field could affect the charge while the charge could never affect the field, there would be no place from which that energy was withdrawn. Action and reaction are not quite the whole argument, but conservation of energy makes the reciprocal coupling unavoidable.

\section{A Changing Source Sends Out a Disturbance}

A point charge produces an electric field directed outward or inward from its position. A current in a wire--a flow of charged particles--produces a magnetic field that circulates around the wire. These familiar static pictures already show two different geometric patterns: electric flux emanates from charge, while magnetic field lines wind around current.

Now move the point charge abruptly. Near its new position the electric field must begin to rearrange, but a distant observer cannot learn of the motion instantly. The old field remains far away until a signal has had time to arrive. Between the near and far regions lies a moving front at which the field is changing. That moving rearrangement is an electromagnetic disturbance.

If the charge is driven back and forth, successive rearrangements follow one another and form an electromagnetic wave. This is the elementary function of an antenna: an oscillating electric current sends a field disturbance outward. Reversing a current in a wire gives the magnetic version of the same argument. The circulation of \(\mathbf B\) must reverse, but the reversal propagates rather than appearing everywhere at once.

This causal picture sets two targets for Maxwell's equations. They must describe how charge and current source the fields, and they must permit changes in those fields to propagate no faster than light. The second target is also a test of relativity. If Maxwell's equations have the same form in every inertial frame and imply a definite wave speed, then every inertial frame must infer the same speed from them.

\section{Reviewing the Relativistic Wave Equation}

Before turning to vector fields, consider a scalar field \(\phi(\mathbf x,t)\), one number at every spacetime point. A first-derivative equation
\begin{equation}
\partial_\mu\phi=0
\end{equation}
would say that every spacetime derivative vanishes. The field would be constant and therefore unable to carry a wave. A nontrivial relativistic wave equation uses two derivatives:
\begin{equation}
\partial_\mu\partial^\mu\phi=0.
\end{equation}
With the sign convention used here,
\begin{equation}
\left(\partial_t^2-\partial_x^2-\partial_y^2-\partial_z^2\right)\phi=0,
\label{eq:scalar-wave-four}
\end{equation}
or, more compactly,
\begin{equation}
\left(\partial_t^2-\nabla^2\right)\phi=0.
\end{equation}
The choice \(c=1\) lets time and distance carry compatible units. Restoring ordinary units would place \(1/c^2\) beside the second time derivative.

For a wave traveling along the \(z\)-axis, the field depends on \(z\) and \(t\), not on \(x\) or \(y\). Equation~\eqref{eq:scalar-wave-four} reduces to
\begin{equation}
\left(\partial_t^2-\partial_z^2\right)\phi=0.
\label{eq:scalar-wave-one}
\end{equation}

\begin{figure}[H]
\centering
\includegraphics[width=0.56\textwidth]{lecture_58_figure_04.png}
\caption{The one-dimensional scalar wave equation used as the model for electromagnetic waves.}
\lectureframe{00:47:03}
\end{figure}

Its general traveling-wave form is
\begin{equation}
\phi(z,t)=f(z-t)+g(z+t).
\end{equation}
The profile \(f(z-t)\) moves toward increasing \(z\), because a point of fixed phase satisfies \(z-t=\text{constant}\). The profile \(g(z+t)\) moves in the opposite direction. Nothing requires either profile to be sinusoidal; any sufficiently smooth shape translates without changing form.

Sinusoids are nevertheless convenient:
\begin{equation}
\phi(z,t)=A\cos\!\bigl(k(z-t)\bigr).
\end{equation}
The wave number \(k\) controls the wavelength,
\begin{equation}
\lambda=\frac{2\pi}{k}.
\end{equation}
A large \(k\) produces rapid spatial oscillation and a short wavelength; a small \(k\) produces a long wavelength. These same profiles will soon become components of \(\mathbf E\) and \(\mathbf B\).

\section{The Vector Calculus Maxwell Needs}

Maxwell's equations are written with dot products, cross products, divergence, and curl. It is worth rebuilding the notation once, because the later component calculations use precisely these patterns.

For ordinary three-vectors \(\mathbf a\) and \(\mathbf b\), the dot product is the scalar
\begin{equation}
\mathbf a\cdot\mathbf b
=a_xb_x+a_yb_y+a_zb_z.
\end{equation}
The cross product is another vector,
\begin{align}
(\mathbf a\times\mathbf b)_x&=a_yb_z-a_zb_y,\\
(\mathbf a\times\mathbf b)_y&=a_zb_x-a_xb_z,\\
(\mathbf a\times\mathbf b)_z&=a_xb_y-a_yb_x.
\end{align}
The remaining components follow from the cyclic order \(x\to y\to z\to x\).

Introduce the differential operator
\begin{equation}
\boldsymbol\nabla=
\left(\partial_x,\partial_y,\partial_z\right).
\end{equation}
Its entries are operators rather than ordinary numbers, but the dot- and cross-product patterns can still be used. For a vector field \(\mathbf A\),
\begin{equation}
\boldsymbol\nabla\cdot\mathbf A
=\partial_xA_x+\partial_yA_y+\partial_zA_z
\label{eq:divergence}
\end{equation}
is the divergence, while
\begin{align}
(\boldsymbol\nabla\times\mathbf A)_x
 &=\partial_yA_z-\partial_zA_y,\\
(\boldsymbol\nabla\times\mathbf A)_y
 &=\partial_zA_x-\partial_xA_z,\\
(\boldsymbol\nabla\times\mathbf A)_z
 &=\partial_xA_y-\partial_yA_x
\label{eq:curl}
\end{align}
are the components of the curl.

The geometry is as important as the formulas. A vector field has positive divergence where it appears to flow outward from a source. It has curl where it circulates. The radial electric field of a point charge illustrates divergence; the magnetic field winding around a current-carrying wire illustrates curl. A spatially constant vector field has neither.

\section{Maxwell's Equations in Vacuum}

In the absence of charge and current, Maxwell's equations are
\begin{align}
\boldsymbol\nabla\times\mathbf E&=-\partial_t\mathbf B,
\label{eq:faraday-vac}\\
\boldsymbol\nabla\cdot\mathbf B&=0,
\label{eq:gauss-b-vac}\\
\boldsymbol\nabla\times\mathbf B&=\partial_t\mathbf E,
\label{eq:ampere-vac}\\
\boldsymbol\nabla\cdot\mathbf E&=0.
\label{eq:gauss-e-vac}
\end{align}
The two curl equations nearly exchange \(\mathbf E\) and \(\mathbf B\), but an essential minus sign distinguishes them. The two divergence equations say that neither field has a source in the vacuum region under consideration. Magnetic divergence remains zero even after electric sources are restored; ordinary electromagnetism contains no magnetic charge.

\begin{figure}[H]
\centering
\includegraphics[width=0.76\textwidth]{lecture_58_figure_06.png}
\caption{The four vacuum Maxwell equations, arranged to display the near symmetry between electric and magnetic fields.}
\lectureframe{00:27:30}
\end{figure}

Charge density \(\rho\) will eventually replace the zero on the right of Equation~\eqref{eq:gauss-e-vac}, and current density \(\mathbf J\) will enter Equation~\eqref{eq:ampere-vac}. We postpone those terms temporarily. The source-free equations expose the transformation structure and the wave mechanism with the least clutter.

They also sharpen Einstein's puzzle. A Galilean observer following a wave would expect to subtract velocities and measure a smaller wave speed. Yet if the same Maxwell equations hold in every inertial frame, each observer derives the same unit wave speed. The way out is not to alter the wave equation from frame to frame. It is to use the Lorentz transformation and to exhibit Maxwell's equations in objects whose transformation laws are already known.

\section{Electric and Magnetic Fields as One Tensor}

The electric and magnetic fields combine into the antisymmetric field tensor
\begin{equation}
F^{\mu\nu}=
\begin{pmatrix}
0   & -E_x & -E_y & -E_z\\
E_x & 0    & -B_z & B_y\\
E_y & B_z  & 0    & -B_x\\
E_z & -B_y & B_x  & 0
\end{pmatrix},
\label{eq:field-tensor}
\end{equation}
where the row and column order is \(t,x,y,z\). Antisymmetry,
\begin{equation}
F^{\mu\nu}=-F^{\nu\mu},
\end{equation}
forces the diagonal entries to vanish and determines the lower half of the matrix from the upper half. The time-space entries contain \(\mathbf E\); the purely spatial entries contain \(\mathbf B\).

\begin{figure}[H]
\centering
\includegraphics[width=0.58\textwidth]{lecture_58_figure_03.png}
\caption{The completed electromagnetic field tensor, including its electric and magnetic entries.}
\lectureframe{00:32:05}
\end{figure}

The importance of \(F^{\mu\nu}\) is not merely that six field components fit into one matrix. Under a Lorentz transformation it transforms as a rank-two tensor. Electric and magnetic fields measured by different observers therefore mix in exactly the way required for the entire object to transform coherently.

There is a second antisymmetric tensor, the dual field tensor \(\widetilde F^{\mu\nu}\). With the present convention it is obtained from Equation~\eqref{eq:field-tensor} by the replacement
\begin{equation}
\mathbf E\longmapsto-\mathbf B,
\qquad
\mathbf B\longmapsto\mathbf E.
\label{eq:dual-replacement}
\end{equation}
Thus
\begin{equation}
\widetilde F^{\mu\nu}=
\begin{pmatrix}
0    & B_x  & B_y  & B_z\\
-B_x & 0    & -E_z & E_y\\
-B_y & E_z  & 0    & -E_x\\
-B_z & -E_y & E_x  & 0
\end{pmatrix}.
\label{eq:dual-tensor}
\end{equation}
The sign in Equation~\eqref{eq:dual-replacement} matters. The lecture first tries the interchange verbally, catches the sign mismatch, and corrects it before completing the matrix.

\begin{figure}[H]
\centering
\includegraphics[width=0.58\textwidth]{lecture_58_figure_07.png}
\caption{The dual tensor being assembled by exchanging electric and magnetic entries with the required sign.}
\lectureframe{00:35:45}
\end{figure}

\section{The Covariant Vacuum Equations}

Each Maxwell equation contains one derivative. Starting from a rank-two tensor and contracting one derivative index leaves one free index, so the natural covariant candidates are
\begin{align}
\partial_\mu F^{\mu\nu}&=0,
\label{eq:maxwell-covariant-vac}\\
\partial_\mu\widetilde F^{\mu\nu}&=0.
\label{eq:maxwell-dual-vac}
\end{align}
For each fixed value of \(\nu\), Equation~\eqref{eq:maxwell-covariant-vac} is one equation. Since \(\nu\) can be \(t,x,y,z\), it packages four scalar equations. Equation~\eqref{eq:maxwell-dual-vac} packages the other four. This matches the two vector curl equations and two scalar divergence equations in the original list.

Compact notation is useful only if it reproduces the familiar equations. Take \(\nu=z\):
\begin{align}
\partial_\mu F^{\mu z}
&=\partial_tF^{tz}
  +\partial_xF^{xz}
  +\partial_yF^{yz}
  +\partial_zF^{zz}\\
&=-\partial_tE_z+\partial_xB_y-\partial_yB_x\\
&=-\partial_tE_z+
  \bigl(\boldsymbol\nabla\times\mathbf B\bigr)_z.
\end{align}
The term \(F^{zz}\) vanishes by antisymmetry. Setting the result to zero gives the \(z\)-component of
\begin{equation}
\boldsymbol\nabla\times\mathbf B=\partial_t\mathbf E.
\end{equation}
The \(x\)- and \(y\)-components follow by cycling the indices.

\begin{figure}[H]
\centering
\includegraphics[width=0.76\textwidth]{lecture_58_figure_08.png}
\caption{Component expansion of the covariant field equation; the diagonal tensor entry vanishes and the remaining terms reconstruct Maxwell's equations.}
\lectureframe{00:43:30}
\end{figure}

Now take the time component:
\begin{align}
\partial_\mu F^{\mu t}
&=\partial_tF^{tt}
  +\partial_xF^{xt}
  +\partial_yF^{yt}
  +\partial_zF^{zt}\\
&=\partial_xE_x+\partial_yE_y+\partial_zE_z\\
&=\boldsymbol\nabla\cdot\mathbf E.
\end{align}
It therefore gives Equation~\eqref{eq:gauss-e-vac}. Performing the same checks on the dual tensor yields
\(\boldsymbol\nabla\cdot\mathbf B=0\) and
\(\boldsymbol\nabla\times\mathbf E=-\partial_t\mathbf B\).

The compact equations are now established rather than merely guessed. Their free-index structure makes them four-vector equations, so a Lorentz transformation changes their components but not their form. This is the precise sense in which the vacuum Maxwell equations are the same in every inertial frame.

\section{From First-Order Equations to a Wave}

Maxwell's equations contain only first derivatives, whereas a wave equation contains second derivatives. The apparent mismatch is resolved because Maxwell's theory couples two fields. A pair of first-order equations for \(\mathbf E\) and \(\mathbf B\) can be combined into a second-order equation for either field alone.

Differentiate Equation~\eqref{eq:ampere-vac} with respect to time:
\begin{equation}
\partial_t^2\mathbf E
=\boldsymbol\nabla\times\partial_t\mathbf B.
\end{equation}
Faraday's law supplies the time derivative of \(\mathbf B\),
\begin{equation}
\partial_t\mathbf B=-\boldsymbol\nabla\times\mathbf E,
\end{equation}
and therefore
\begin{equation}
\partial_t^2\mathbf E
=-\boldsymbol\nabla\times
  \bigl(\boldsymbol\nabla\times\mathbf E\bigr).
\label{eq:double-curl}
\end{equation}

The board works through one component rather than invoking a memorized identity. For the \(x\)-component,
\begin{align}
\bigl[\boldsymbol\nabla\times
(\boldsymbol\nabla\times\mathbf E)\bigr]_x
&=\partial_y(\boldsymbol\nabla\times\mathbf E)_z
 -\partial_z(\boldsymbol\nabla\times\mathbf E)_y\\
&=\partial_y(\partial_xE_y-\partial_yE_x)
 -\partial_z(\partial_zE_x-\partial_xE_z)\\
&=\partial_x(\partial_yE_y+\partial_zE_z)
 -\partial_y^2E_x-\partial_z^2E_x.
\end{align}
Add and subtract \(\partial_x^2E_x\) to expose the divergence:
\begin{align}
\bigl[\boldsymbol\nabla\times
(\boldsymbol\nabla\times\mathbf E)\bigr]_x
&=\partial_x
  (\partial_xE_x+\partial_yE_y+\partial_zE_z)
 -(\partial_x^2+\partial_y^2+\partial_z^2)E_x\\
&=\partial_x(\boldsymbol\nabla\cdot\mathbf E)
 -\nabla^2E_x.
\label{eq:double-curl-x}
\end{align}
In vacuum \(\boldsymbol\nabla\cdot\mathbf E=0\). Combining
Equations~\eqref{eq:double-curl} and \eqref{eq:double-curl-x} gives
\begin{equation}
\left(\partial_t^2-\nabla^2\right)E_x=0.
\label{eq:ex-wave}
\end{equation}

\begin{figure}[H]
\centering
\includegraphics[width=0.72\textwidth]{lecture_58_figure_09.png}
\caption{The component calculation reaches the wave equation for \(E_x\).}
\lectureframe{00:57:10}
\end{figure}

The same calculation applies to \(E_y\) and \(E_z\), and an analogous elimination gives
\begin{equation}
\left(\partial_t^2-\nabla^2\right)\mathbf B=0.
\end{equation}
Electric and magnetic fields can therefore propagate through empty space as waves. In the units used here their speed is one, which means the speed is \(c\) when ordinary units are restored.

\section{A Plane Wave and Its Polarization}

Construct a simple wave traveling in the positive \(z\)-direction:
\begin{equation}
\mathbf E(z,t)
=\widehat{\mathbf x}\,E_0
  \cos\!\bigl(k(z-t)\bigr).
\label{eq:plane-e}
\end{equation}
Its profile translates along \(z\) without changing shape. Because the electric field points along the \(x\)-axis, the wave is linearly polarized in the \(x\)-direction. We have chosen \(E_y=0\) only to keep the example simple; a second transverse polarization could be included.

The longitudinal component is not equally free. Since the wave depends only on \(z\) and \(t\),
\begin{equation}
\boldsymbol\nabla\cdot\mathbf E
=\partial_zE_z=0.
\end{equation}
Thus \(E_z\) may be a constant, but it cannot contain the varying wave profile. A uniform background field can be superposed or removed independently; it is not part of the propagating radiation. The wave itself is transverse:
\begin{equation}
\mathbf E_{\mathrm{wave}}\cdot\widehat{\mathbf z}=0.
\end{equation}

Faraday's law fixes the magnetic field. For Equation~\eqref{eq:plane-e}, the only nonzero component of \(\boldsymbol\nabla\times\mathbf E\) is the \(y\)-component,
\begin{equation}
(\boldsymbol\nabla\times\mathbf E)_y=\partial_zE_x.
\end{equation}
Consequently \(\partial_t\mathbf B\) points along \(y\), and a compatible choice is
\begin{equation}
\mathbf B(z,t)
=\widehat{\mathbf y}\,E_0
  \cos\!\bigl(k(z-t)\bigr).
\label{eq:plane-b}
\end{equation}
In these units the electric and magnetic amplitudes are equal. Their orientation obeys
\begin{equation}
\mathbf E\perp\mathbf B,
\qquad
\mathbf E\perp\widehat{\mathbf z},
\qquad
\mathbf B\perp\widehat{\mathbf z},
\qquad
\mathbf E\times\mathbf B\parallel\widehat{\mathbf z}.
\end{equation}

\begin{figure}[H]
\centering
\includegraphics[width=0.70\textwidth]{lecture_58_figure_05.png}
\caption{The completed board sketch of a transverse electromagnetic wave: the electric and magnetic oscillations are perpendicular to one another and to the direction of propagation.}
\lectureframe{01:05:25}
\end{figure}

\begin{figure}[H]
\centering
\begin{tikzpicture}[x=1cm,y=1cm,>=Latex]
  \draw[->,thick] (0,0) -- (8.2,0) node[right] {$z$};
  \draw[->,thick] (0,-1.6) -- (0,1.7) node[above] {$x$};
  \draw[->,thick] (0,0) -- (-1.15,-0.8) node[below left] {$y$};
  \draw[blue!70!black,thick,domain=0:7.5,samples=120]
    plot (\x,{0.95*sin(90*\x)});
  \foreach \z/\s in {0.8/1,2.8/-1,4.8/1,6.8/-1}{
    \draw[->,green!45!black,thick]
      (\z,0) -- ++({-0.55*\s},{-0.38*\s});
  }
  \node[blue!70!black] at (6.9,1.25) {$\mathbf E$};
  \node[green!45!black] at (5.2,-1.05) {$\mathbf B$};
  \draw[->,very thick] (6.2,1.55) -- (7.7,1.55)
    node[midway,above] {propagation};
\end{tikzpicture}
\caption{A clean geometric reconstruction: \(\mathbf E\), \(\mathbf B\), and the propagation direction form a mutually perpendicular triad.}
\end{figure}

This is the character of light and of all electromagnetic radiation in vacuum. Since the equations that produced the wave are Lorentz covariant, another inertial observer again describes an electromagnetic wave moving at \(c\), with transverse electric and magnetic fields.

\section{Charge Density and Current Density}

We now restore the sources that were temporarily set aside. When many charged particles occupy a region, it is useful to coarse-grain them into a continuous charge density
\begin{equation}
\rho(\mathbf x,t)=\frac{dQ}{dV}.
\end{equation}
The approximation is like treating water as a continuous fluid even though it consists of molecules. Positive and negative charges contribute with opposite signs, so \(\rho\) is the net charge per unit volume.

The charge in a chosen region \(V\) is
\begin{equation}
Q_V(t)=\int_V\rho(\mathbf x,t)\,dV.
\label{eq:charge-volume}
\end{equation}
The density may vary with position and time. A pile of charge can move away from a point, changing the local density there. It cannot, however, vary in an arbitrary fashion. Nature conserves electric charge.

Global conservation alone would be too weak. It would allow charge to disappear from Earth and reappear on Mars at the same instant while the total remained unchanged. The stronger statement is local: charge can leave one place and reach another only through a flow in the intervening space.

To describe that flow, take an infinitesimal oriented surface element
\begin{equation}
d\boldsymbol\sigma=\widehat{\mathbf n}\,d\sigma.
\end{equation}
Its magnitude is the area, and its direction is normal to the surface. The current density \(\mathbf J\) is defined so that the charge crossing the element per unit time is
\begin{equation}
\frac{dQ}{dt}=\mathbf J\cdot d\boldsymbol\sigma.
\label{eq:current-flux-element}
\end{equation}
The magnitude of \(\mathbf J\) is charge per unit area per unit time, and its direction is the direction of positive-charge flow.

\begin{figure}[H]
\centering
\includegraphics[width=0.64\textwidth]{lecture_58_figure_10.png}
\caption{An oriented area element and the flux factor \(d\boldsymbol\sigma\cdot\mathbf J\).}
\lectureframe{01:18:15}
\end{figure}

The dot product selects the component normal to the surface. A current tangent to the surface crosses neither side and contributes zero. A current parallel to the normal gives the maximum flux for a fixed magnitude. A wire current is one special case, but \(\mathbf J\) may also describe charge flowing through a sheet or throughout three-dimensional space.

\section{Local Conservation and the Continuity Equation}

Fix a region \(V\) in space and orient every element of its closed boundary \(\partial V\) outward. If there is a net outward current, the charge inside decreases:
\begin{equation}
\frac{dQ_V}{dt}
=-\oint_{\partial V}\mathbf J\cdot d\boldsymbol\sigma.
\label{eq:integral-conservation}
\end{equation}
The minus sign is physical rather than conventional: positive outward flux removes positive charge from the region. Using Equation~\eqref{eq:charge-volume},
\begin{equation}
\int_V\partial_t\rho\,dV
=-\oint_{\partial V}\mathbf J\cdot d\boldsymbol\sigma.
\label{eq:integral-continuity-one}
\end{equation}

\begin{figure}[H]
\centering
\includegraphics[width=0.72\textwidth]{lecture_58_figure_11.png}
\caption{The fixed-region conservation law after the surface flux is converted into a volume statement.}
\lectureframe{01:29:00}
\end{figure}

Gauss's theorem, also called the divergence theorem, relates flux through a closed surface to divergence in its interior:
\begin{equation}
\oint_{\partial V}\mathbf A\cdot d\boldsymbol\sigma
=\int_V\boldsymbol\nabla\cdot\mathbf A\,dV.
\label{eq:divergence-theorem}
\end{equation}
One can picture hidden pipes injecting water at points inside a vessel. The divergence measures the local rate at which water emerges from the pipe ends; the surface integral measures the total water leaving the vessel. Whatever is created as outward flow inside must eventually cross the boundary.

Applying Equation~\eqref{eq:divergence-theorem} to \(\mathbf J\) gives
\begin{equation}
\int_V
\left(\partial_t\rho+\boldsymbol\nabla\cdot\mathbf J\right)dV=0.
\end{equation}
This relation holds for every possible region \(V\). The integrand must therefore vanish point by point:
\begin{equation}
\boxed{\partial_t\rho+\boldsymbol\nabla\cdot\mathbf J=0.}
\label{eq:continuity-three}
\end{equation}

\begin{figure}[H]
\centering
\includegraphics[width=0.72\textwidth]{lecture_58_figure_12.png}
\caption{The local continuity equation written component by component.}
\lectureframe{01:29:50}
\end{figure}

Equation~\eqref{eq:continuity-three} says exactly what local conservation requires. If charge density decreases at a point or within a small region, current must diverge outward from it. If current converges inward, charge density increases.

The density and current naturally form the four-current
\begin{equation}
J^\mu=(\rho,J_x,J_y,J_z)
\label{eq:four-current}
\end{equation}
in \(c=1\) units. The continuity equation becomes the Lorentz scalar equation
\begin{equation}
\partial_\mu J^\mu=0.
\label{eq:continuity-four}
\end{equation}
The flow of a relativistic distribution should transform as a four-vector, just as the four-velocity of one particle does. Equation~\eqref{eq:continuity-four} makes local charge conservation valid in every inertial frame.

\section{Maxwell's Equations with Sources}

Restoring charge and current gives
\begin{align}
\boldsymbol\nabla\cdot\mathbf E&=\rho,
\label{eq:gauss-e-source}\\
\boldsymbol\nabla\cdot\mathbf B&=0,\\
\boldsymbol\nabla\times\mathbf E&=-\partial_t\mathbf B,\\
\boldsymbol\nabla\times\mathbf B&=\partial_t\mathbf E+\mathbf J.
\label{eq:ampere-source}
\end{align}
The equations involving electric sources collapse into
\begin{equation}
\boxed{\partial_\mu F^{\mu\nu}=J^\nu,}
\label{eq:maxwell-source-covariant}
\end{equation}
while the source-free magnetic pair remains
\begin{equation}
\boxed{\partial_\mu\widetilde F^{\mu\nu}=0.}
\label{eq:maxwell-dual-source}
\end{equation}

\begin{figure}[H]
\centering
\includegraphics[width=0.78\textwidth]{lecture_58_figure_13.png}
\caption{The sourced Maxwell equations beside their covariant form \(\partial_\mu F^{\mu\nu}=J^\nu\).}
\lectureframe{01:34:00}
\end{figure}

The time component of Equation~\eqref{eq:maxwell-source-covariant} reads
\begin{equation}
\partial_\mu F^{\mu t}=J^t=\rho,
\end{equation}
and the component calculation already performed turns its left side into
\(\boldsymbol\nabla\cdot\mathbf E\). The three spatial components reproduce
\(\boldsymbol\nabla\times\mathbf B=\partial_t\mathbf E+\mathbf J\).

There is also a useful consistency check. Applying \(\partial_\nu\) to
Equation~\eqref{eq:maxwell-source-covariant} gives
\begin{equation}
\partial_\nu J^\nu
=\partial_\nu\partial_\mu F^{\mu\nu}=0,
\end{equation}
because the two derivatives are symmetric under \(\mu\leftrightarrow\nu\), whereas \(F^{\mu\nu}\) is antisymmetric.\editorialnote{This short consistency check follows from the covariant equation and makes explicit why Maxwell's sourced equations automatically respect the continuity equation.}

The reciprocal structure is now complete. The tensor force law
\(f^\mu=qF^\mu{}_\nu u^\nu\) describes how fields act on charge, and
Equation~\eqref{eq:maxwell-source-covariant} describes how charge and current generate fields. Both are tensor equations. In vacuum, the same field equations yield transverse waves moving at \(c\). Lorentz covariance and the electromagnetic wave equation therefore meet at the point that motivated the construction: every inertial observer obtains the same laws of electromagnetism and the same speed of light.
